% ============================================================
% Quantum Mamba architecture diagram (TikZ)
% Include in main document with:
%   \usepackage{tikz}
%   \usetikzlibrary{positioning, arrows.meta, fit, backgrounds, calc, shapes.geometric}
%   % ============================================================
% Quantum Mamba architecture diagram (TikZ)
% Include in main document with:
%   \usepackage{tikz}
%   \usetikzlibrary{positioning, arrows.meta, fit, backgrounds, calc, shapes.geometric}
%   % ============================================================
% Quantum Mamba architecture diagram (TikZ)
% Include in main document with:
%   \usepackage{tikz}
%   \usetikzlibrary{positioning, arrows.meta, fit, backgrounds, calc, shapes.geometric}
%   % ============================================================
% Quantum Mamba architecture diagram (TikZ)
% Include in main document with:
%   \usepackage{tikz}
%   \usetikzlibrary{positioning, arrows.meta, fit, backgrounds, calc, shapes.geometric}
%   \input{fig_quantum_mamba.tex}
% ============================================================

\begin{figure*}[t]
\centering
\begin{tikzpicture}[
    font=\small,
    node distance=4mm and 6mm,
    >=Latex,
    classical/.style={
        rectangle, rounded corners=2pt,
        draw=black!70, thick,
        fill=blue!8,
        minimum width=28mm, minimum height=7mm,
        align=center
    },
    cbranch/.style={
        rectangle, rounded corners=2pt,
        draw=black!70, thick,
        fill=blue!8,
        minimum width=22mm, minimum height=6mm,
        align=center, font=\footnotesize
    },
    qbranch/.style={
        rectangle, rounded corners=2pt,
        draw=purple!80!black, thick,
        fill=purple!8,
        minimum width=32mm, minimum height=9mm,
        align=center, font=\footnotesize
    },
    qcore/.style={
        rectangle, rounded corners=2pt,
        draw=purple!80!black, thick,
        fill=purple!6,
        minimum width=28mm, minimum height=7mm,
        align=center, font=\footnotesize
    },
    tensor/.style={font=\footnotesize\itshape},
    arrow/.style={-{Latex[length=2mm]}, thick},
]

% ---------- Input ----------
\node[classical, fill=gray!10] (input) {Input sequence\\$\mathbf{X} \in \mathbb{R}^{B \times T \times F}$};

% ---------- Classical encoder ----------
\node[classical, below=of input] (encoder) {Classical encoder\\\footnotesize Linear $+$ LN $+$ SiLU};

% ---------- Chunking ----------
\node[classical, below=of encoder] (chunk) {Chunking $+$ intra-chunk attention pooling};
\node[tensor, right=2mm of chunk] {$\{\mathbf{s}^{(k)}\}_{k=1}^{K}$};

% ---------- C→Q projections ----------
\node[classical, below=of chunk, minimum width=70mm] (proj) {%
    Classical$\to$quantum parameter projections\\%
    \footnotesize $\boldsymbol{\phi}^{(k)},\;\boldsymbol{\theta}_{\text{fwd}}^{(k)},\;\boldsymbol{\theta}_{\text{bwd}}^{(k)},\;\boldsymbol{\theta}_{\text{diag}}^{(k)},\;\Delta^{(k)}$
};

% ---------- Quantum core (big box) ----------
% Title above the branches
\node[below=6mm of proj, font=\footnotesize\bfseries, text=purple!70!black] (coretitle)
    {Three-branch coherent quantum core};

% Input encoding
\node[qcore, below=2mm of coretitle] (enc) {%
    $|\psi_0\rangle = \bigotimes_i R_Y(\phi_i^{(k)})|0\rangle$
};

% Three branches side by side
\node[qbranch, below left=6mm and 6mm of enc] (fwd) {%
    \textbf{Forward branch}\\%
    $U_{\text{fwd}}(\boldsymbol{\theta}_{\text{fwd}}^{(k)},\Delta^{(k)})$\\%
    $\Rightarrow |\psi_1^{(k)}\rangle$
};
\node[qbranch, below=6mm of enc] (bwd) {%
    \textbf{Backward branch}\\%
    $U_{\text{bwd}}(\boldsymbol{\theta}_{\text{bwd}}^{(k)},\Delta^{(k)})$\\%
    $\Rightarrow |\psi_2^{(k)}\rangle$
};
\node[qbranch, below right=6mm and 6mm of enc] (diag) {%
    \textbf{Diagonal branch}\\%
    $U_{\text{diag}}(\boldsymbol{\theta}_{\text{diag}}^{(k)})$\\%
    $\Rightarrow |\psi_3^{(k)}\rangle$
};

% Coherent superposition
\node[qcore, below=6mm of bwd, minimum width=75mm] (super) {%
    Coherent superposition:\; $|\psi^{(k)}\rangle = \mathcal{N}\!\left(\alpha|\psi_1^{(k)}\rangle + \beta|\psi_2^{(k)}\rangle + \gamma|\psi_3^{(k)}\rangle\right)$
};

% Measurement
\node[qcore, below=3mm of super, minimum width=75mm] (meas) {%
    Multi-observable measurement\; $\langle\sigma_X^{(i)}\rangle,\langle\sigma_Y^{(i)}\rangle,\langle\sigma_Z^{(i)}\rangle \Rightarrow \mathbf{q}^{(k)} \in \mathbb{R}^{3n}$
};

% Bounding box for the core
\begin{pgfonlayer}{background}
\node[
    draw=purple!60!black, thick, dashed, rounded corners=3pt,
    fill=purple!3,
    fit=(coretitle)(enc)(fwd)(diag)(super)(meas),
    inner sep=3mm
] (qcoreBox) {};
\end{pgfonlayer}

% ---------- Output projection ----------
\node[classical, below=4mm of qcoreBox] (outproj) {%
    Output projection:\; $\mathbf{c}^{(k)} = \mathrm{Linear}(\mathbf{q}^{(k)})$
};

% ---------- Sequence aggregation ----------
\node[classical, below=of outproj] (seqagg) {Sequence aggregation $\Rightarrow \mathbf{r} \in \mathbb{R}^{d}$};

% ---------- Classifier ----------
\node[classical, below=of seqagg, fill=green!8] (cls) {Classification head $\Rightarrow$ logits};

% ---------- Arrows ----------
\draw[arrow] (input) -- (encoder);
\draw[arrow] (encoder) -- (chunk);
\draw[arrow] (chunk) -- (proj);
\draw[arrow] (proj) -- (qcoreBox.north);

\draw[arrow] (enc) -- (fwd.north);
\draw[arrow] (enc) -- (bwd);
\draw[arrow] (enc) -- (diag.north);

\draw[arrow] (fwd.south) -- ([xshift=-20mm]super.north);
\draw[arrow] (bwd.south) -- (super.north);
\draw[arrow] (diag.south) -- ([xshift=20mm]super.north);

\draw[arrow] (super) -- (meas);
\draw[arrow] (qcoreBox.south) -- (outproj);
\draw[arrow] (outproj) -- (seqagg);
\draw[arrow] (seqagg) -- (cls);

\end{tikzpicture}
\caption{%
\textbf{Quantum Mamba architecture.}
A classical encoder projects the input sequence $\mathbf{X}$ and partitions it into chunks, each summarized by intra-chunk attention pooling.
For each chunk, the summary $\mathbf{s}^{(k)}$ is mapped through separate linear projections to the quantum input angles $\boldsymbol{\phi}^{(k)}$, three sets of variational parameters $\{\boldsymbol{\theta}_{\text{fwd}}^{(k)}, \boldsymbol{\theta}_{\text{bwd}}^{(k)}, \boldsymbol{\theta}_{\text{diag}}^{(k)}\}$, and a scalar selective-gating factor $\Delta^{(k)}$.
Three independent quantum branches (forward and backward Delta-modulated, diagonal unmodulated) are evolved from the same input state $|\psi_0\rangle$ and coherently combined at the amplitude level through learnable complex coefficients $\alpha, \beta, \gamma$ before multi-observable measurement.
The resulting chunk-level quantum feature vectors $\mathbf{q}^{(k)}$ are then projected back to the model hidden space, aggregated across chunks, and passed through a classification head.%
}
\label{fig:qmamba_architecture}
\end{figure*}

% ============================================================

\begin{figure*}[t]
\centering
\begin{tikzpicture}[
    font=\small,
    node distance=4mm and 6mm,
    >=Latex,
    classical/.style={
        rectangle, rounded corners=2pt,
        draw=black!70, thick,
        fill=blue!8,
        minimum width=28mm, minimum height=7mm,
        align=center
    },
    cbranch/.style={
        rectangle, rounded corners=2pt,
        draw=black!70, thick,
        fill=blue!8,
        minimum width=22mm, minimum height=6mm,
        align=center, font=\footnotesize
    },
    qbranch/.style={
        rectangle, rounded corners=2pt,
        draw=purple!80!black, thick,
        fill=purple!8,
        minimum width=32mm, minimum height=9mm,
        align=center, font=\footnotesize
    },
    qcore/.style={
        rectangle, rounded corners=2pt,
        draw=purple!80!black, thick,
        fill=purple!6,
        minimum width=28mm, minimum height=7mm,
        align=center, font=\footnotesize
    },
    tensor/.style={font=\footnotesize\itshape},
    arrow/.style={-{Latex[length=2mm]}, thick},
]

% ---------- Input ----------
\node[classical, fill=gray!10] (input) {Input sequence\\$\mathbf{X} \in \mathbb{R}^{B \times T \times F}$};

% ---------- Classical encoder ----------
\node[classical, below=of input] (encoder) {Classical encoder\\\footnotesize Linear $+$ LN $+$ SiLU};

% ---------- Chunking ----------
\node[classical, below=of encoder] (chunk) {Chunking $+$ intra-chunk attention pooling};
\node[tensor, right=2mm of chunk] {$\{\mathbf{s}^{(k)}\}_{k=1}^{K}$};

% ---------- C→Q projections ----------
\node[classical, below=of chunk, minimum width=70mm] (proj) {%
    Classical$\to$quantum parameter projections\\%
    \footnotesize $\boldsymbol{\phi}^{(k)},\;\boldsymbol{\theta}_{\text{fwd}}^{(k)},\;\boldsymbol{\theta}_{\text{bwd}}^{(k)},\;\boldsymbol{\theta}_{\text{diag}}^{(k)},\;\Delta^{(k)}$
};

% ---------- Quantum core (big box) ----------
% Title above the branches
\node[below=6mm of proj, font=\footnotesize\bfseries, text=purple!70!black] (coretitle)
    {Three-branch coherent quantum core};

% Input encoding
\node[qcore, below=2mm of coretitle] (enc) {%
    $|\psi_0\rangle = \bigotimes_i R_Y(\phi_i^{(k)})|0\rangle$
};

% Three branches side by side
\node[qbranch, below left=6mm and 6mm of enc] (fwd) {%
    \textbf{Forward branch}\\%
    $U_{\text{fwd}}(\boldsymbol{\theta}_{\text{fwd}}^{(k)},\Delta^{(k)})$\\%
    $\Rightarrow |\psi_1^{(k)}\rangle$
};
\node[qbranch, below=6mm of enc] (bwd) {%
    \textbf{Backward branch}\\%
    $U_{\text{bwd}}(\boldsymbol{\theta}_{\text{bwd}}^{(k)},\Delta^{(k)})$\\%
    $\Rightarrow |\psi_2^{(k)}\rangle$
};
\node[qbranch, below right=6mm and 6mm of enc] (diag) {%
    \textbf{Diagonal branch}\\%
    $U_{\text{diag}}(\boldsymbol{\theta}_{\text{diag}}^{(k)})$\\%
    $\Rightarrow |\psi_3^{(k)}\rangle$
};

% Coherent superposition
\node[qcore, below=6mm of bwd, minimum width=75mm] (super) {%
    Coherent superposition:\; $|\psi^{(k)}\rangle = \mathcal{N}\!\left(\alpha|\psi_1^{(k)}\rangle + \beta|\psi_2^{(k)}\rangle + \gamma|\psi_3^{(k)}\rangle\right)$
};

% Measurement
\node[qcore, below=3mm of super, minimum width=75mm] (meas) {%
    Multi-observable measurement\; $\langle\sigma_X^{(i)}\rangle,\langle\sigma_Y^{(i)}\rangle,\langle\sigma_Z^{(i)}\rangle \Rightarrow \mathbf{q}^{(k)} \in \mathbb{R}^{3n}$
};

% Bounding box for the core
\begin{pgfonlayer}{background}
\node[
    draw=purple!60!black, thick, dashed, rounded corners=3pt,
    fill=purple!3,
    fit=(coretitle)(enc)(fwd)(diag)(super)(meas),
    inner sep=3mm
] (qcoreBox) {};
\end{pgfonlayer}

% ---------- Output projection ----------
\node[classical, below=4mm of qcoreBox] (outproj) {%
    Output projection:\; $\mathbf{c}^{(k)} = \mathrm{Linear}(\mathbf{q}^{(k)})$
};

% ---------- Sequence aggregation ----------
\node[classical, below=of outproj] (seqagg) {Sequence aggregation $\Rightarrow \mathbf{r} \in \mathbb{R}^{d}$};

% ---------- Classifier ----------
\node[classical, below=of seqagg, fill=green!8] (cls) {Classification head $\Rightarrow$ logits};

% ---------- Arrows ----------
\draw[arrow] (input) -- (encoder);
\draw[arrow] (encoder) -- (chunk);
\draw[arrow] (chunk) -- (proj);
\draw[arrow] (proj) -- (qcoreBox.north);

\draw[arrow] (enc) -- (fwd.north);
\draw[arrow] (enc) -- (bwd);
\draw[arrow] (enc) -- (diag.north);

\draw[arrow] (fwd.south) -- ([xshift=-20mm]super.north);
\draw[arrow] (bwd.south) -- (super.north);
\draw[arrow] (diag.south) -- ([xshift=20mm]super.north);

\draw[arrow] (super) -- (meas);
\draw[arrow] (qcoreBox.south) -- (outproj);
\draw[arrow] (outproj) -- (seqagg);
\draw[arrow] (seqagg) -- (cls);

\end{tikzpicture}
\caption{%
\textbf{Quantum Mamba architecture.}
A classical encoder projects the input sequence $\mathbf{X}$ and partitions it into chunks, each summarized by intra-chunk attention pooling.
For each chunk, the summary $\mathbf{s}^{(k)}$ is mapped through separate linear projections to the quantum input angles $\boldsymbol{\phi}^{(k)}$, three sets of variational parameters $\{\boldsymbol{\theta}_{\text{fwd}}^{(k)}, \boldsymbol{\theta}_{\text{bwd}}^{(k)}, \boldsymbol{\theta}_{\text{diag}}^{(k)}\}$, and a scalar selective-gating factor $\Delta^{(k)}$.
Three independent quantum branches (forward and backward Delta-modulated, diagonal unmodulated) are evolved from the same input state $|\psi_0\rangle$ and coherently combined at the amplitude level through learnable complex coefficients $\alpha, \beta, \gamma$ before multi-observable measurement.
The resulting chunk-level quantum feature vectors $\mathbf{q}^{(k)}$ are then projected back to the model hidden space, aggregated across chunks, and passed through a classification head.%
}
\label{fig:qmamba_architecture}
\end{figure*}

% ============================================================

\begin{figure*}[t]
\centering
\begin{tikzpicture}[
    font=\small,
    node distance=4mm and 6mm,
    >=Latex,
    classical/.style={
        rectangle, rounded corners=2pt,
        draw=black!70, thick,
        fill=blue!8,
        minimum width=28mm, minimum height=7mm,
        align=center
    },
    cbranch/.style={
        rectangle, rounded corners=2pt,
        draw=black!70, thick,
        fill=blue!8,
        minimum width=22mm, minimum height=6mm,
        align=center, font=\footnotesize
    },
    qbranch/.style={
        rectangle, rounded corners=2pt,
        draw=purple!80!black, thick,
        fill=purple!8,
        minimum width=32mm, minimum height=9mm,
        align=center, font=\footnotesize
    },
    qcore/.style={
        rectangle, rounded corners=2pt,
        draw=purple!80!black, thick,
        fill=purple!6,
        minimum width=28mm, minimum height=7mm,
        align=center, font=\footnotesize
    },
    tensor/.style={font=\footnotesize\itshape},
    arrow/.style={-{Latex[length=2mm]}, thick},
]

% ---------- Input ----------
\node[classical, fill=gray!10] (input) {Input sequence\\$\mathbf{X} \in \mathbb{R}^{B \times T \times F}$};

% ---------- Classical encoder ----------
\node[classical, below=of input] (encoder) {Classical encoder\\\footnotesize Linear $+$ LN $+$ SiLU};

% ---------- Chunking ----------
\node[classical, below=of encoder] (chunk) {Chunking $+$ intra-chunk attention pooling};
\node[tensor, right=2mm of chunk] {$\{\mathbf{s}^{(k)}\}_{k=1}^{K}$};

% ---------- C→Q projections ----------
\node[classical, below=of chunk, minimum width=70mm] (proj) {%
    Classical$\to$quantum parameter projections\\%
    \footnotesize $\boldsymbol{\phi}^{(k)},\;\boldsymbol{\theta}_{\text{fwd}}^{(k)},\;\boldsymbol{\theta}_{\text{bwd}}^{(k)},\;\boldsymbol{\theta}_{\text{diag}}^{(k)},\;\Delta^{(k)}$
};

% ---------- Quantum core (big box) ----------
% Title above the branches
\node[below=6mm of proj, font=\footnotesize\bfseries, text=purple!70!black] (coretitle)
    {Three-branch coherent quantum core};

% Input encoding
\node[qcore, below=2mm of coretitle] (enc) {%
    $|\psi_0\rangle = \bigotimes_i R_Y(\phi_i^{(k)})|0\rangle$
};

% Three branches side by side
\node[qbranch, below left=6mm and 6mm of enc] (fwd) {%
    \textbf{Forward branch}\\%
    $U_{\text{fwd}}(\boldsymbol{\theta}_{\text{fwd}}^{(k)},\Delta^{(k)})$\\%
    $\Rightarrow |\psi_1^{(k)}\rangle$
};
\node[qbranch, below=6mm of enc] (bwd) {%
    \textbf{Backward branch}\\%
    $U_{\text{bwd}}(\boldsymbol{\theta}_{\text{bwd}}^{(k)},\Delta^{(k)})$\\%
    $\Rightarrow |\psi_2^{(k)}\rangle$
};
\node[qbranch, below right=6mm and 6mm of enc] (diag) {%
    \textbf{Diagonal branch}\\%
    $U_{\text{diag}}(\boldsymbol{\theta}_{\text{diag}}^{(k)})$\\%
    $\Rightarrow |\psi_3^{(k)}\rangle$
};

% Coherent superposition
\node[qcore, below=6mm of bwd, minimum width=75mm] (super) {%
    Coherent superposition:\; $|\psi^{(k)}\rangle = \mathcal{N}\!\left(\alpha|\psi_1^{(k)}\rangle + \beta|\psi_2^{(k)}\rangle + \gamma|\psi_3^{(k)}\rangle\right)$
};

% Measurement
\node[qcore, below=3mm of super, minimum width=75mm] (meas) {%
    Multi-observable measurement\; $\langle\sigma_X^{(i)}\rangle,\langle\sigma_Y^{(i)}\rangle,\langle\sigma_Z^{(i)}\rangle \Rightarrow \mathbf{q}^{(k)} \in \mathbb{R}^{3n}$
};

% Bounding box for the core
\begin{pgfonlayer}{background}
\node[
    draw=purple!60!black, thick, dashed, rounded corners=3pt,
    fill=purple!3,
    fit=(coretitle)(enc)(fwd)(diag)(super)(meas),
    inner sep=3mm
] (qcoreBox) {};
\end{pgfonlayer}

% ---------- Output projection ----------
\node[classical, below=4mm of qcoreBox] (outproj) {%
    Output projection:\; $\mathbf{c}^{(k)} = \mathrm{Linear}(\mathbf{q}^{(k)})$
};

% ---------- Sequence aggregation ----------
\node[classical, below=of outproj] (seqagg) {Sequence aggregation $\Rightarrow \mathbf{r} \in \mathbb{R}^{d}$};

% ---------- Classifier ----------
\node[classical, below=of seqagg, fill=green!8] (cls) {Classification head $\Rightarrow$ logits};

% ---------- Arrows ----------
\draw[arrow] (input) -- (encoder);
\draw[arrow] (encoder) -- (chunk);
\draw[arrow] (chunk) -- (proj);
\draw[arrow] (proj) -- (qcoreBox.north);

\draw[arrow] (enc) -- (fwd.north);
\draw[arrow] (enc) -- (bwd);
\draw[arrow] (enc) -- (diag.north);

\draw[arrow] (fwd.south) -- ([xshift=-20mm]super.north);
\draw[arrow] (bwd.south) -- (super.north);
\draw[arrow] (diag.south) -- ([xshift=20mm]super.north);

\draw[arrow] (super) -- (meas);
\draw[arrow] (qcoreBox.south) -- (outproj);
\draw[arrow] (outproj) -- (seqagg);
\draw[arrow] (seqagg) -- (cls);

\end{tikzpicture}
\caption{%
\textbf{Quantum Mamba architecture.}
A classical encoder projects the input sequence $\mathbf{X}$ and partitions it into chunks, each summarized by intra-chunk attention pooling.
For each chunk, the summary $\mathbf{s}^{(k)}$ is mapped through separate linear projections to the quantum input angles $\boldsymbol{\phi}^{(k)}$, three sets of variational parameters $\{\boldsymbol{\theta}_{\text{fwd}}^{(k)}, \boldsymbol{\theta}_{\text{bwd}}^{(k)}, \boldsymbol{\theta}_{\text{diag}}^{(k)}\}$, and a scalar selective-gating factor $\Delta^{(k)}$.
Three independent quantum branches (forward and backward Delta-modulated, diagonal unmodulated) are evolved from the same input state $|\psi_0\rangle$ and coherently combined at the amplitude level through learnable complex coefficients $\alpha, \beta, \gamma$ before multi-observable measurement.
The resulting chunk-level quantum feature vectors $\mathbf{q}^{(k)}$ are then projected back to the model hidden space, aggregated across chunks, and passed through a classification head.%
}
\label{fig:qmamba_architecture}
\end{figure*}

% ============================================================

\begin{figure*}[t]
\centering
\begin{tikzpicture}[
    font=\small,
    node distance=4mm and 6mm,
    >=Latex,
    classical/.style={
        rectangle, rounded corners=2pt,
        draw=black!70, thick,
        fill=blue!8,
        minimum width=28mm, minimum height=7mm,
        align=center
    },
    cbranch/.style={
        rectangle, rounded corners=2pt,
        draw=black!70, thick,
        fill=blue!8,
        minimum width=22mm, minimum height=6mm,
        align=center, font=\footnotesize
    },
    qbranch/.style={
        rectangle, rounded corners=2pt,
        draw=purple!80!black, thick,
        fill=purple!8,
        minimum width=32mm, minimum height=9mm,
        align=center, font=\footnotesize
    },
    qcore/.style={
        rectangle, rounded corners=2pt,
        draw=purple!80!black, thick,
        fill=purple!6,
        minimum width=28mm, minimum height=7mm,
        align=center, font=\footnotesize
    },
    tensor/.style={font=\footnotesize\itshape},
    arrow/.style={-{Latex[length=2mm]}, thick},
]

% ---------- Input ----------
\node[classical, fill=gray!10] (input) {Input sequence\\$\mathbf{X} \in \mathbb{R}^{B \times T \times F}$};

% ---------- Classical encoder ----------
\node[classical, below=of input] (encoder) {Classical encoder\\\footnotesize Linear $+$ LN $+$ SiLU};

% ---------- Chunking ----------
\node[classical, below=of encoder] (chunk) {Chunking $+$ intra-chunk attention pooling};
\node[tensor, right=2mm of chunk] {$\{\mathbf{s}^{(k)}\}_{k=1}^{K}$};

% ---------- C→Q projections ----------
\node[classical, below=of chunk, minimum width=70mm] (proj) {%
    Classical$\to$quantum parameter projections\\%
    \footnotesize $\boldsymbol{\phi}^{(k)},\;\boldsymbol{\theta}_{\text{fwd}}^{(k)},\;\boldsymbol{\theta}_{\text{bwd}}^{(k)},\;\boldsymbol{\theta}_{\text{diag}}^{(k)},\;\Delta^{(k)}$
};

% ---------- Quantum core (big box) ----------
% Title above the branches
\node[below=6mm of proj, font=\footnotesize\bfseries, text=purple!70!black] (coretitle)
    {Three-branch coherent quantum core};

% Input encoding
\node[qcore, below=2mm of coretitle] (enc) {%
    $|\psi_0\rangle = \bigotimes_i R_Y(\phi_i^{(k)})|0\rangle$
};

% Three branches side by side
\node[qbranch, below left=6mm and 6mm of enc] (fwd) {%
    \textbf{Forward branch}\\%
    $U_{\text{fwd}}(\boldsymbol{\theta}_{\text{fwd}}^{(k)},\Delta^{(k)})$\\%
    $\Rightarrow |\psi_1^{(k)}\rangle$
};
\node[qbranch, below=6mm of enc] (bwd) {%
    \textbf{Backward branch}\\%
    $U_{\text{bwd}}(\boldsymbol{\theta}_{\text{bwd}}^{(k)},\Delta^{(k)})$\\%
    $\Rightarrow |\psi_2^{(k)}\rangle$
};
\node[qbranch, below right=6mm and 6mm of enc] (diag) {%
    \textbf{Diagonal branch}\\%
    $U_{\text{diag}}(\boldsymbol{\theta}_{\text{diag}}^{(k)})$\\%
    $\Rightarrow |\psi_3^{(k)}\rangle$
};

% Coherent superposition
\node[qcore, below=6mm of bwd, minimum width=75mm] (super) {%
    Coherent superposition:\; $|\psi^{(k)}\rangle = \mathcal{N}\!\left(\alpha|\psi_1^{(k)}\rangle + \beta|\psi_2^{(k)}\rangle + \gamma|\psi_3^{(k)}\rangle\right)$
};

% Measurement
\node[qcore, below=3mm of super, minimum width=75mm] (meas) {%
    Multi-observable measurement\; $\langle\sigma_X^{(i)}\rangle,\langle\sigma_Y^{(i)}\rangle,\langle\sigma_Z^{(i)}\rangle \Rightarrow \mathbf{q}^{(k)} \in \mathbb{R}^{3n}$
};

% Bounding box for the core
\begin{pgfonlayer}{background}
\node[
    draw=purple!60!black, thick, dashed, rounded corners=3pt,
    fill=purple!3,
    fit=(coretitle)(enc)(fwd)(diag)(super)(meas),
    inner sep=3mm
] (qcoreBox) {};
\end{pgfonlayer}

% ---------- Output projection ----------
\node[classical, below=4mm of qcoreBox] (outproj) {%
    Output projection:\; $\mathbf{c}^{(k)} = \mathrm{Linear}(\mathbf{q}^{(k)})$
};

% ---------- Sequence aggregation ----------
\node[classical, below=of outproj] (seqagg) {Sequence aggregation $\Rightarrow \mathbf{r} \in \mathbb{R}^{d}$};

% ---------- Classifier ----------
\node[classical, below=of seqagg, fill=green!8] (cls) {Classification head $\Rightarrow$ logits};

% ---------- Arrows ----------
\draw[arrow] (input) -- (encoder);
\draw[arrow] (encoder) -- (chunk);
\draw[arrow] (chunk) -- (proj);
\draw[arrow] (proj) -- (qcoreBox.north);

\draw[arrow] (enc) -- (fwd.north);
\draw[arrow] (enc) -- (bwd);
\draw[arrow] (enc) -- (diag.north);

\draw[arrow] (fwd.south) -- ([xshift=-20mm]super.north);
\draw[arrow] (bwd.south) -- (super.north);
\draw[arrow] (diag.south) -- ([xshift=20mm]super.north);

\draw[arrow] (super) -- (meas);
\draw[arrow] (qcoreBox.south) -- (outproj);
\draw[arrow] (outproj) -- (seqagg);
\draw[arrow] (seqagg) -- (cls);

\end{tikzpicture}
\caption{%
\textbf{Quantum Mamba architecture.}
A classical encoder projects the input sequence $\mathbf{X}$ and partitions it into chunks, each summarized by intra-chunk attention pooling.
For each chunk, the summary $\mathbf{s}^{(k)}$ is mapped through separate linear projections to the quantum input angles $\boldsymbol{\phi}^{(k)}$, three sets of variational parameters $\{\boldsymbol{\theta}_{\text{fwd}}^{(k)}, \boldsymbol{\theta}_{\text{bwd}}^{(k)}, \boldsymbol{\theta}_{\text{diag}}^{(k)}\}$, and a scalar selective-gating factor $\Delta^{(k)}$.
Three independent quantum branches (forward and backward Delta-modulated, diagonal unmodulated) are evolved from the same input state $|\psi_0\rangle$ and coherently combined at the amplitude level through learnable complex coefficients $\alpha, \beta, \gamma$ before multi-observable measurement.
The resulting chunk-level quantum feature vectors $\mathbf{q}^{(k)}$ are then projected back to the model hidden space, aggregated across chunks, and passed through a classification head.%
}
\label{fig:qmamba_architecture}
\end{figure*}
